\chapter{Where the Crypto Venues Live}\label{ch:nw:where-the-crypto-venues-live}

Many of the largest crypto venues run their matching engines in or next to a \omterm{def:nw:trading-in-a-public-cloud:cloud}{public cloud} region rather than in an exchange's own data
centre, and for several of them the cloud provider says so. An AWS blog of June 2026 puts Binance in Tokyo, Bybit in Singapore, Coinbase in the United States and Deribit in London;
research reported by CoinDesk in March 2026 puts Hyperliquid's validators in AWS's Tokyo region, alongside Binance and KuCoin infrastructure;
BitMEX moved its platform from AWS's Irish region to Tokyo because more than 200 milliseconds separated it from the firms that traded it. A
firm that wants to be close launches machines in the venue's region, measures, and keeps the few that happen to sit nearest the venue's
servers.

Chapter~16 described the cloud; this chapter is about finding a venue in it. What is publicly known about where the venues are; how to find
out when it is not published (the provider's address ranges, the round trips from known places); the lottery of launching machines until one
lands close; and the re-measurement that notices when the venue moves. Everything that is not a venue's or a provider's own statement is a
labelled simulation.

\section{The publicly known locations}

\begin{dated}{2026-09}{Where the cloud-hosted crypto venues are, from public sources}\label{dat:nw:where-the-crypto-venues-live:venues}
\begin{itemize}
\item AWS (June 2026): ``Major exchanges operate across Regions: Binance in Tokyo, Bybit in Singapore, Coinbase in the US, and Deribit in London.''
\item Glassnode research reported by CoinDesk (March 2026): Hyperliquid's 24 validators across several \omterm{def:nw:trading-in-a-public-cloud:azid}{availability zones} of AWS's
  ap-northeast-1 (Tokyo) region, its API behind AWS CloudFront; Hyperliquid, Binance and KuCoin concentrating infrastructure in ap-northeast-1;
  Tokyo users reaching the validators in 2 to 3 milliseconds, European users in more than 200.
\item AWS case study: BitMEX moved its trading platform from the Europe (Ireland) Region to Asia Pacific (Tokyo), with a ten-minute cutover,
  order execution below 10 milliseconds and liquidity up more than 185\% across core contracts.
\end{itemize}
\end{dated}

The box is short because the sources are few. Venues rarely publish their region in their own documentation; the providers mention them in
blogs and case studies; researchers infer them from measurements. A region is also only a start: chapter~16 showed that inside it the zone
matters by hundreds of microseconds, and the venue's zone identifier is published even more rarely.

\begin{definition}[Autonomous system, anycast, edge network]\label{def:nw:where-the-crypto-venues-live:as}
An \emph{autonomous system}\index{autonomous system} (AS) is a group of IP prefixes run by one or more network operators under a single routing
policy, identified by a number in the internet's routing system. \emph{Anycast}\index{anycast} is the practice of making one service address
available in several places at once, so that each packet reaches whichever is nearest in the routing system. An \emph{edge
network}\index{edge network} is a provider's set of points of presence close to users, often reached by anycast, that terminates connections
and forwards requests to the origin servers.
\end{definition}

\omterm{def:nw:where-the-crypto-venues-live:as}{Anycast} and \omterm{def:nw:where-the-crypto-venues-live:as}{edge networks} are why looking up a venue's public address often says nothing about its matching engine. Hyperliquid's API, in the
report above, is reached through the provider's \omterm{def:nw:where-the-crypto-venues-live:as}{edge network}, while its validators sit in one region; a user in Europe resolves the same name
as a user in Tokyo, reaches a nearby edge, and still waits for the round trip to Tokyo behind it.

\section{How to discover a location}

A provider that publishes its address ranges gives the first clue. AWS publishes a file, ip-ranges.json, listing each prefix with its region and
the service that uses it. Matching a venue's endpoint addresses against it by longest prefix says which region an address belongs to, or that
it belongs to the \omterm{def:nw:where-the-crypto-venues-live:as}{edge network}, in which case it says nothing about the origin.

\omcode{code/firm/venuefind/firm_venuefind.py}{38}{53}{A published IP-range file, loaded and matched by longest prefix.}\label{lst:nw:where-the-crypto-venues-live:match}

\begin{table}[ht]
\centering\small
\begin{tabular}{llll}
\toprule
Endpoint address & Longest matching prefix & Region & Service \\
\midrule
192.0.2.200 & 192.0.2.128/25 & ap-northeast-1 & EC2 \\
198.51.100.7 & 198.51.100.0/24 & ap-southeast-1 & EC2 \\
203.0.113.9 & 203.0.113.0/24 & GLOBAL & CLOUDFRONT \\
2001:db8:1abc::1 & 2001:db8:1000::/36 & ap-northeast-1 & EC2 \\
\bottomrule
\end{tabular}
\caption{Four endpoint addresses matched against the chapter's fixture of a published IP-range file (the documented format; addresses from the
ranges reserved for documentation, so no real host is named). The third address is on the \omterm{def:nw:where-the-crypto-venues-live:as}{edge network}: the lookup locates the edge, not the
venue. Data: \texttt{nw\_find.lookups()}.}\label{tab:nw:where-the-crypto-venues-live:lookups}
\end{table}

\begin{definition}[Round-trip triangulation]\label{def:nw:where-the-crypto-venues-live:tri}
\emph{Round-trip triangulation}\index{round-trip triangulation} estimates where a server is from round trips measured to it from probes at known
places: each round trip bounds the distance (light cannot cover more than half of it in glass) and, divided by a \omterm{def:nw:the-north-american-map:factor}{route factor}, estimates it; the
estimate is the point whose distances best agree with those estimates.
\end{definition}

\begin{proposition}[The light-time bound]\label{prop:nw:where-the-crypto-venues-live:bound}
A server that answers a probe in a round trip $T$, of which $E$ is equipment and processing on the way, lies within $(T/2 - E)\,c_0/n_g$ of the
probe along the geodesic, for the smallest \omterm{def:nw:the-north-american-map:floor}{group index} $n_g$ on the path.
\end{proposition}

\begin{proof}
Half the round trip less the equipment is the most time the signal can have spent travelling one way; at speed $c_0/n_g$ it covers at most that
distance, and no path is shorter than the geodesic.
\end{proof}

The bound is exact; the estimate is not. In the chapter's example, probes in Hong Kong, Singapore, Sydney and Mumbai measure round trips of
36.5, 67.4, 98.8 and \qty{85.5}{\milli\second} to a server hidden in Tokyo (a labelled simulation, with paths 30\% longer than the geodesics
and some noise). Triangulation with the right \omterm{def:nw:the-north-american-map:factor}{route factor}, 1.3, puts the server within \qty{23}{\kilo\metre} of the truth; with 1.2 or 1.4 it
misses by more than \qty{540}{\kilo\metre}. Triangulation from other cities finds the region; inside the region the firm must measure from
inside.

\begin{omfigure}
\begin{tikzpicture}
\begin{axis}[axis equal image,width=11cm,xmin=-6200,xmax=6600,ymin=-6200,ymax=3200,xlabel={km east of 120\textdegree{} E},
  ylabel={km north of 15\textdegree{} N},tick label style={font=\footnotesize,/pgf/number format/1000 sep={\,}},label style={font=\small},grid=major,grid style={gray!20},
  xtick distance=2000,table/col sep=comma,scaled ticks=false]
\addplot[only marks,mark=*,mark size=2.2pt,omDef,nodes near coords,point meta=explicit symbolic,
  every node near coord/.append style={font=\scriptsize,anchor=west}] table[x=x,y=y,meta=name,restrict expr to domain={\coordindex}{0:3}]{figdata/networks/17-where-the-crypto-venues-live/triangulation.csv};
\addplot[only marks,mark=square*,mark size=2.4pt,omThm] table[x=x,y=y,restrict expr to domain={\coordindex}{4:4}]{figdata/networks/17-where-the-crypto-venues-live/triangulation.csv};
\addplot[only marks,mark=x,mark size=4pt,very thick,omMeth] table[x=x,y=y,restrict expr to domain={\coordindex}{5:5}]{figdata/networks/17-where-the-crypto-venues-live/triangulation.csv};
\node[font=\scriptsize,anchor=south east,omThm] at (axis cs:2100,2450) {server (hidden) and estimate};
\end{axis}
\end{tikzpicture}

\omcaption{Simulation: \omterm{def:nw:where-the-crypto-venues-live:tri}{round-trip triangulation} from four probes in other cities to a server in Tokyo, with the right \omterm{def:nw:the-north-american-map:factor}{route factor}; the estimate
(cross) falls \qty{23}{\kilo\metre} from the hidden server (square). A local projection for drawing only. Data: \texttt{fig\_find.py},
\texttt{nw\_find.triangulate\_example()}.}\label{fig:nw:where-the-crypto-venues-live:tri}
\end{omfigure}

\omcode{code/firm/venuefind/firm_venuefind.py}{60}{81}{Round-trip triangulation: a grid search for the point whose distances best match the round
trips, never beyond any probe's light-time bound.}\label{lst:nw:where-the-crypto-venues-live:tri}

\section{The instance lottery}

\begin{definition}[Instance lottery]\label{def:nw:where-the-crypto-venues-live:lottery}
The \emph{instance lottery}\index{instance lottery} is the practice of launching many cloud machines in a venue's region and zone, measuring each
one's round trip to the venue, keeping the fastest few and releasing the rest, because the provider places each machine without telling the
customer where.
\end{definition}

Inside one zone, where a new machine lands decides its distance from the venue's servers: the same rack, the same row, another hall. The
provider does not say, so the firm draws. AWS's blog of June 2026 reports that \omterm{def:nw:trading-in-a-public-cloud:pg}{placement groups}, adapter tuning and kernel bypass ``can achieve
sub-150~µs intra-Region latencies'', and chapter~16's measurement puts a typical round trip within a zone at \qty{270}{\micro\second}. The model
takes those two figures as the median and the best 1\% of launches, with lognormal round trips between them: an assumption, labelled.

\begin{proposition}[Best of $n$]\label{prop:nw:where-the-crypto-venues-live:bestn}
If launches' round trips are independent draws from a distribution $F$, the expected best of $n$ is $\int_0^\infty (1 - F(t))^n\,dt$, which falls
with $n$ ever more slowly. If each launch costs $c$ and each microsecond saved is worth $v$, the best number of launches maximises
$v\,(E[X_{(1:1)}] - E[X_{(1:n)}]) - c\,n$.
\end{proposition}

\begin{proof}
The minimum of $n$ independent draws exceeds $t$ only if all $n$ do, with probability $(1-F(t))^n$; the expectation of a positive variable is the
integral of its survival function. The net value is the saving against a single launch less the launches' cost.
\end{proof}

\begin{omfigure}
\begin{tikzpicture}
\begin{semilogxaxis}[width=11cm,height=5.6cm,xmin=1,xmax=160,ymin=120,ymax=290,xlabel={instances launched (log)},
  ylabel={expected best round trip (\si{\micro\second})},tick label style={font=\footnotesize},label style={font=\small},grid=major,
  grid style={gray!25},table/col sep=comma]
\addplot[omDef,thick,mark=*,mark size=1.3pt] table[x=n,y=best_us]{figdata/networks/17-where-the-crypto-venues-live/lottery.csv};
\end{semilogxaxis}
\end{tikzpicture}

\omcaption{Simulation: the expected best round trip to the venue among $n$ launches, for lognormal launches with a median of
\qty{270}{\micro\second} and a best 1\% at \qty{150}{\micro\second} (the published figures of chapter~16 and of the AWS blog, used as assumptions).
The first ten launches save \qty{94}{\micro\second}; the next 150 another 47. Data: \texttt{fig\_find.py}, \texttt{nw\_find.lottery\_curve()}.}\label{fig:nw:where-the-crypto-venues-live:lottery}
\end{omfigure}

The curve is the order statistics of \cref{prop:nw:where-the-crypto-venues-live:bestn}: one launch expects \qty{279}{\micro\second}, ten
185, forty 157, a hundred and sixty \qty{138}{\micro\second}. Whether the next launch is worth it depends on what a microsecond is
worth. At an assumed 100~USD a month per microsecond and a launch that costs a day of a 16-vCPU instance (17.14~USD at chapter~16's price), the
best is 77 launches, worth about 11\,700~USD a month net of their cost (\cref{fig:nw:where-the-crypto-venues-live:net}).

\omcode{code/firm/venuefind/firm_venuefind.py}{90}{108}{The instance lottery: the expected best of $n$ launches by simulation, and the number of
launches that maximises the value saved net of their cost.}\label{lst:nw:where-the-crypto-venues-live:lottery}

\begin{omfigure}
\begin{tikzpicture}
\begin{axis}[width=11cm,height=5.4cm,xmin=0,xmax=150,ymin=-0.5,ymax=13,xlabel={instances launched},
  ylabel={net value (thousand USD a month)},tick label style={font=\footnotesize},label style={font=\small},grid=major,
  grid style={gray!25},table/col sep=comma]
\addplot[omThm,thick,mark=*,mark size=1.3pt] table[x=n,y=net_k]{figdata/networks/17-where-the-crypto-venues-live/net.csv};
\draw[omMeth,dashed] (axis cs:77,-0.5) -- (axis cs:77,13) node[pos=0.1,right,font=\scriptsize]{best: 77};
\end{axis}
\end{tikzpicture}

\omcaption{Simulation: the monthly value of the round trip saved against a single launch, net of the launches' cost, for an assumed value of 100~USD a
month per microsecond and 17.14~USD per launch; the optimum is flat, and anything from 40 to 120 launches is within a few per cent of it. Data:
\texttt{fig\_find.py}.}\label{fig:nw:where-the-crypto-venues-live:net}
\end{omfigure}

\section{Continuous re-measurement}

The lottery's winners do not stay winners. The venue may move within the region or to another region, as BitMEX did; the provider may migrate
machines; maintenance may reshuffle a zone. A firm that measured once will notice only when its fills get worse. The chapter's check is Book~7's
CUSUM test on the daily median round trip from each kept machine: the cumulative sum of standardised recursive residuals, flagged when it leaves
its 5\% boundary.

\begin{omfigure}
\begin{tikzpicture}
\begin{axis}[width=11cm,height=5.4cm,xmin=0,xmax=91,ymin=140,ymax=205,xlabel={day},ylabel={daily median round trip (\si{\micro\second})},
  tick label style={font=\footnotesize},label style={font=\small},grid=major,grid style={gray!25},table/col sep=comma]
\addplot[omDef,thick,no markers] table[x=day,y=median_us]{figdata/networks/17-where-the-crypto-venues-live/drift.csv};
\draw[gray,dashed] (axis cs:61,140) -- (axis cs:61,205) node[pos=0.92,left,font=\scriptsize]{move};
\draw[omMeth,thick] (axis cs:73,140) -- (axis cs:73,205) node[pos=0.92,right,font=\scriptsize]{flag};
\end{axis}
\end{tikzpicture}

\omcaption{Simulation: a machine's daily median round trip to a venue rises from 160 to \qty{185}{\micro\second} when the venue moves on day 61;
Book~7's CUSUM test flags the change on day 73. Data: \texttt{fig\_find.py}, \texttt{nw\_find.drift\_example()}.}\label{fig:nw:where-the-crypto-venues-live:drift}
\end{omfigure}

The CUSUM test is slow by design: it was built to detect a change in the mean without crying wolf, and here it needs twelve days to be sure of a
shift of four standard deviations. A firm that races will add a faster alarm (a jump in the median beyond a fixed threshold, confirmed the next day)
and use the CUSUM test to confirm; the lottery then starts again.

\begin{method}[Finding and following a cloud-hosted venue]\label{met:nw:where-the-crypto-venues-live:find}
\begin{enumerate}
\item Read what is public: the venue's documentation and announcements, the provider's blogs and case studies, research; record each claim with its
  source and date.
\item Match the venue's endpoint addresses against the provider's published ranges; distinguish origin regions from \omterm{def:nw:where-the-crypto-venues-live:as}{edge networks}.
\item Triangulate from probes in other regions to confirm the region; inside it, launch machines in every zone and measure.
\item Run the lottery in the best zone: launch, measure with hardware timestamps, keep the best, release the rest; stop where the next launch is worth
  less than it costs.
\item Re-measure every day; flag changes; start again when the venue or the provider moves.
\end{enumerate}
\end{method}

\section{Tutorial: find the venue, win the lottery}\label{tut:nw:where-the-crypto-venues-live:tut}

\begin{tutorial}
\textbf{Goal.} Match addresses, triangulate a server, simulate the lottery and a move, with no cloud account.
\textbf{End state:} \cref{fig:nw:where-the-crypto-venues-live:tri,fig:nw:where-the-crypto-venues-live:lottery,fig:nw:where-the-crypto-venues-live:net,fig:nw:where-the-crypto-venues-live:drift}
and \cref{tab:nw:where-the-crypto-venues-live:lookups}.
\begin{tutsteps}
\item \textbf{Addresses.} \texttt{firm\_venuefind.load\_ranges} reads the fixture (documented format, documentation addresses); \texttt{match}
  (\cref{lst:nw:where-the-crypto-venues-live:match}) finds the longest prefix.
\item \textbf{Triangulation.} \texttt{nw\_find.probes()} simulates round trips from four sites of the site table; \texttt{triangulate}
  (\cref{lst:nw:where-the-crypto-venues-live:tri}) estimates the server.
\item \textbf{Lottery.} \texttt{expected\_best} and \texttt{best\_n} (\cref{lst:nw:where-the-crypto-venues-live:lottery}) with \texttt{LOTTERY},
  \texttt{VALUE} and \texttt{LAUNCH}.
\item \textbf{Re-measurement.} \texttt{moved} wraps \texttt{firm\_decay.cusum}.
\end{tutsteps}
\textbf{What to change next.} Replace the simulated probes with round trips measured from machines you run; give the lottery a two-tier
distribution (same rack or not) and see how the best $n$ changes.
\end{tutorial}

\section{Build: the venue finder}\label{bld:nw:where-the-crypto-venues-live:venuefind}

\begin{build}
\bfield{Purpose} Where a cloud-hosted venue is and how close the firm's machines are to it: the input of chapters~18 and~19 and of the cloud rows of
chapter~29's plan.
\bfield{Interface} \texttt{firm\_venuefind}: \texttt{load\_ranges}, \texttt{match}, \texttt{distance\_bound\_km}, \texttt{triangulate},
\texttt{Lottery}, \texttt{expected\_best}, \texttt{best\_n}, \texttt{moved} (Book~7's \texttt{firm.decay} wrapped, not edited).
\bfield{Rules} Public location claims carry their source and date; an edge address is never taken for an origin; triangulation respects every
light-time bound; the lottery's distribution is a stated assumption until replaced by measurements.
\bfield{Acceptance tests} \texttt{code/firm/venuefind/tests/}: longest-prefix matching (IPv4 and IPv6), triangulation of a known point and refusal of an
impossible one, the lottery's monotonicity, and the CUSUM flag with and without a move.
\bfield{Stretch} Probe round trips through the edge to separate its delay from the origin's; a two-sided change detector tuned for speed.
\end{build}

\begin{omsources}
\item AWS for Industries blog, ``Ultra-low-latency cross-Region crypto trading with Avelacom and AWS'' (June 2026); AWS case study on BitMEX; AWS
documentation of its IP address ranges.
\item CoinDesk, ``Hyperliquid traders in Tokyo get 200-millisecond edge, Glassnode research shows'' (30 March 2026).
\item RFC 1930 (autonomous systems), RFC 4786 (anycast), RFC 5737 and RFC 3849 (documentation address ranges).
\end{omsources}

\section{Exercises}

\begin{exercise}[$\star$]\label{exo:nw:where-the-crypto-venues-live:1}
A probe measures a round trip of \qty{1}{\milli\second} to a server. How far away can the server be, at most, in fibre?
\end{exercise}

\begin{exercise}[$\star$]\label{exo:nw:where-the-crypto-venues-live:2}
Why does matching a venue's API address against a provider's ranges sometimes locate nothing useful?
\end{exercise}

\begin{exercise}[$\star$]\label{exo:nw:where-the-crypto-venues-live:3}
Using \cref{fig:nw:where-the-crypto-venues-live:lottery}, how much does the tenth launch save compared with the first, and the hundred-and-sixtieth
compared with the fortieth?
\end{exercise}

\begin{exercise}[$\star\star$]\label{exo:nw:where-the-crypto-venues-live:4}
Why does triangulation from other cities find a region but not a zone?
\end{exercise}

\begin{exercise}[$\star\star$]\label{exo:nw:where-the-crypto-venues-live:5}
BitMEX moved from Ireland to Tokyo. What changed for a market maker running in Ireland, and what should it have done?
\end{exercise}

\begin{exercise}[$\star\star$]\label{exo:nw:where-the-crypto-venues-live:6}
Why is the CUSUM test slow, and what would you pair it with?
\end{exercise}

\begin{exercise}[$\star\star\star$]\label{exo:nw:where-the-crypto-venues-live:7}
\emph{Coding.} With \texttt{firm.venuefind}, find the best number of launches if a microsecond is worth only 20~USD a month.
\end{exercise}

\begin{exercise}[$\star\star\star$]\label{exo:nw:where-the-crypto-venues-live:8}
\emph{Find the flaw.} ``Our machine pings the venue's API in \qty{1}{\milli\second} from Frankfurt, so the venue is in Frankfurt.''
\end{exercise}

\section{Problem: Launch Forty, Keep One}

\begin{problem}[{Weekend problem --- how many machines to launch}]\label{pb:nw:where-the-crypto-venues-live:1}
A firm trades a venue that runs in one zone of a cloud region. Its launches' round trips to the venue are lognormal with a median of
\qty{270}{\micro\second} and a best 1\% at \qty{150}{\micro\second}. A microsecond is worth 100~USD a month to it; a launch, measured for a day, costs
17.14~USD.

\medskip\noindent\textbf{Part I --- The draws.}
\begin{enumerate}
\item What is the expected round trip of a single launch?
\item And the expected best of 10, 40 and 160?
\item Why is the expected single round trip above the median?
\item What does \cref{prop:nw:where-the-crypto-venues-live:bestn} say about the shape of the curve?
\end{enumerate}

\medskip\noindent\textbf{Part II --- The economics.}
\begin{enumerate}[resume]
\item What is the net value of keeping the best of 40?
\item What number of launches maximises the net value, and what is it worth?
\item How flat is the optimum?
\item What if a microsecond is worth 20~USD a month?
\end{enumerate}

\medskip\noindent\textbf{Part III --- Where the draws come from.}
\begin{enumerate}[resume]
\item Which facts of chapter~16 and of the AWS blog does the model use, and which assumptions?
\item Why may launches not be independent?
\item What else than the instance's position changes its round trip?
\item How would you replace the assumption with data?
\end{enumerate}

\medskip\noindent\textbf{Part IV --- The verdict.}
\begin{enumerate}[resume]
\item State the \emph{named result}: the expected best round trip among $n$ instances and the number of launches that maximises value per microsecond net
  of launch cost.
\item How long does the winner stay a winner?
\item What does a venue that publishes its zone identifier change?
\item Is the lottery fair to the venue's other users?
\item What would a venue do to end it?
\item What does the lottery cost a firm that does not race?
\item Where does this go in the connectivity plan of chapter~29?
\item In one sentence: why launch forty and keep one?
\end{enumerate}
\end{problem}

\section{Interview questions}

\begin{interviewq}[$\star$ \iqroles{developer}]\label{iq:nw:where-the-crypto-venues-live:1}
How would you find out which cloud region a crypto exchange runs in?
\end{interviewq}

\begin{interviewq}[$\star\star$ \iqroles{developer}]\label{iq:nw:where-the-crypto-venues-live:2}
What is \omterm{def:nw:where-the-crypto-venues-live:as}{anycast}, and why can it mislead a latency measurement?
\end{interviewq}

\begin{interviewq}[$\star\star$ \iqroles{developer, researcher}]\label{iq:nw:where-the-crypto-venues-live:3}
You launch 50 instances and measure their round trips to a venue. How many do you keep, and how do you decide?
\end{interviewq}

\begin{interviewq}[$\star\star$ \iqroles{developer}]\label{iq:nw:where-the-crypto-venues-live:4}
How do you detect that a venue has moved its servers?
\end{interviewq}

\begin{interviewq}[$\star\star$ \iqroles{trader}]\label{iq:nw:where-the-crypto-venues-live:5}
Why did so many crypto venues and trading firms end up in Tokyo?
\end{interviewq}

\begin{interviewq}[$\star\star\star$ \iqroles{developer, researcher}]\label{iq:nw:where-the-crypto-venues-live:6}
Estimate the location of a server from round trips measured in four cities. What limits the accuracy?
\end{interviewq}
